\section{Experiments\label{sec:experiments}}

We show that adapters achieve parameter efficient transfer for text tasks.
On the GLUE benchmark~\citep{wang2018glue},
adapter tuning is within $0.4\%$ of full fine-tuning of BERT, but it adds only $3\%$ of the number of parameters trained by fine-tuning.
We confirm this result on a further $17$ public classification tasks and SQuAD question answering.
Analysis shows that adapter-based tuning automatically focuses on the higher layers of the network.

\subsection{Experimental Settings}

We use the public, pre-trained BERT Transformer network as our base model.
To perform classification with BERT, we follow the approach in~\citet{devlin2018bert}.
The first token in each sequence is a special ``classification token''.
We attach a linear layer to the embedding of this token to predict the class label.

Our training procedure also follows~\citet{devlin2018bert}.
We optimize using Adam~\citep{kingma2014adam},
whose learning rate is increased linearly over the first $10\%$ of the steps, and then decayed linearly to zero.
All runs are trained on $4$ Google Cloud TPUs with a batch size of $32$.
For each dataset and algorithm, we run a hyperparameter sweep and select the best model according to accuracy on the validation set.
For the GLUE tasks, we report the test metrics provided by the submission website\footnote{\url{https://gluebenchmark.com/}}.
For the other classification tasks we report test-set accuracy.

We compare to fine-tuning, the current standard for transfer of large pre-trained models,
and the strategy successfully used by BERT.
For $N$ tasks, full fine-tuning requires $N{\times}$ the number of parameters of the pre-trained model.
Our goal is to attain performance equal to fine-tuning, but with fewer total parameters, ideally near to $1{\times}$.

\subsection{GLUE benchmark\label{sec:glue}}

\begin{table*}[t]
\centering
\begin{adjustbox}{max width=\textwidth}
\begin{tabular}{l|ll|rrrrrrrrr|r}
\toprule
{} & \pbox{3cm}{Total num\\ params} & \pbox{3cm}{Trained \\ params / task} & CoLA &     SST &    MRPC &   STS-B &     QQP & MNLI\textsubscript{m} & MNLI\textsubscript{mm} &    QNLI &     RTE &   Total \\
\midrule
BERT\textsubscript{LARGE} &  $9.0\times$ &   $100\%$ &  $60.5$ &  $94.9$ &  $89.3$ &  $87.6$ &  $72.1$ &  $86.7$ &  $85.9$ &  $91.1$ &  $70.1$ &  $80.4$ \\
Adapters ($8$-$256$)      &  $1.3\times$ &   $3.6\%$ &  $59.5$ &  $94.0$ &  $89.5$ &  $86.9$ &  $71.8$ &  $84.9$ &  $85.1$ &  $90.7$ &  $71.5$ &  $80.0$ \\
Adapters ($64$)           &  $1.2\times$ &   $2.1\%$ &  $56.9$ &  $94.2$ &  $89.6$ &  $87.3$ &  $71.8$ &  $85.3$ &  $84.6$ &  $91.4$ &  $68.8$ &  $79.6$ \\
\bottomrule
\end{tabular}
\end{adjustbox}
\caption{
Results on GLUE test sets scored using the GLUE evaluation server.
MRPC and QQP are evaluated using F1 score.
STS-B is evaluated using Spearman's correlation coefficient.
CoLA is evaluated using Matthew's Correlation.
The other tasks are evaluated using accuracy.
Adapter tuning achieves comparable overall score ($80.0$) to full fine-tuning ($80.4$) using $1.3\times$ parameters in total, compared to $9\times$.
Fixing the adapter size to $64$ leads to a slightly decreased overall score of $79.6$ and slightly smaller model.
\label{tab:glue}}
\end{table*}

We first evaluate on GLUE.\footnote{
We omit WNLI as in~\citet{devlin2018bert} because the no current algorithm beats the baseline of predicting the majority class.}
For these datasets, we transfer from the pre-trained BERT\textsubscript{LARGE} model,
which contains $24$ layers, and a total of $330$M parameters, see~\citet{devlin2018bert} for details.
We perform a small hyperparameter sweep for adapter tuning:
We sweep learning rates in $\{3 \cdot 10^{-5}, 3 \cdot 10^{-4}, 3 \cdot 10^{-3}\}$, and number of epochs in $\{3, 20\}$.
We test both using a fixed adapter size (number of units in the bottleneck),
and selecting the best size per task from $\{8, 64, 256\}$.
The adapter size is the only adapter-specific hyperparameter that we tune.
Finally, due to training instability, we re-run $5$ times with different random seeds and select the best model on the validation set.

Table~\ref{tab:glue} summarizes the results.
Adapters achieve a mean GLUE score of $80.0$, compared to $80.4$ achieved by full fine-tuning.
The optimal adapter size varies per dataset. For example, $256$ is chosen for MNLI, whereas for the smallest dataset, RTE, $8$ is chosen.
Restricting always to size $64$, leads to a small decrease in average accuracy to $79.6$.
To solve all of the datasets in Table~\ref{tab:glue}, fine-tuning requires $9\times$ the total number of BERT parameters.\footnote{
We treat MNLI\textsubscript{m} and MNLI\textsubscript{mm} as separate tasks with individually tuned hyperparameters.
However, they could be combined into one model, leaving $8\times$ overall.}
In contrast, adapters require only $1.3\times$ parameters.

\subsection{Additional Classification Tasks}

\renewcommand{\arraystretch}{0.9}
\begin{table*}[t]
\centering
\begin{adjustbox}{max width=\textwidth}
\begin{tabular}{l|rrrrr}
\toprule
Dataset &
\pbox{5cm}{No BERT\\baseline} &
\pbox{5cm}{BERT\textsubscript{BASE}\\Fine-tune} &
\pbox{5cm}{BERT\textsubscript{BASE}\\Variable FT} &
\pbox{5cm}{BERT\textsubscript{BASE}\\Adapters} \\
\midrule
20 newsgroups & $ 91.1 $ & $ 92.8 \pm 0.1 $ & $ 92.8 \pm 0.1 $ & $91.7 \pm 0.2$ \\
Crowdflower airline & $ 84.5 $ & $ 83.6 \pm 0.3 $ & $ 84.0 \pm 0.1 $ & $ 84.5 \pm 0.2 $ \\
Crowdflower corporate messaging & $ 91.9 $ & $ 92.5 \pm 0.5 $ & $ 92.4 \pm 0.6 $ & $ 92.9 \pm 0.3 $ \\
Crowdflower disasters & $ 84.9 $ & $ 85.3 \pm 0.4 $ & $ 85.3 \pm 0.4 $ & $ 84.1 \pm 0.2 $ \\
Crowdflower economic news relevance & $ 81.1 $ & $ 82.1 \pm 0.0 $ & $ 78.9 \pm 2.8 $ & $ 82.5 \pm 0.3 $ \\
Crowdflower emotion & $ 36.3 $ & $ 38.4 \pm 0.1 $ & $ 37.6 \pm 0.2 $ & $ 38.7 \pm 0.1 $ \\
Crowdflower global warming & $ 82.7 $ & $ 84.2 \pm 0.4
 $ & $ 81.9 \pm 0.2 $ & $ 82.7 \pm 0.3 $ \\
Crowdflower political audience & $ 81.0 $ & $ 80.9 \pm 0.3 $ & $ 80.7 \pm 0.8
 $ & $ 79.0 \pm 0.5 $ \\
Crowdflower political bias & $ 76.8 $ & $ 75.2 \pm 0.9 $ & $ 76.5 \pm 0.4 $ & $ 75.9 \pm 0.3 $ \\
Crowdflower political message & $ 43.8 $ & $ 38.9 \pm 0.6 $ & $ 44.9 \pm 0.6 $ & $ 44.1 \pm 0.2 $ \\
Crowdflower primary emotions & $ 33.5 $ & $ 36.9 \pm 1.6 $ & $ 38.2 \pm 1.0 $ & $ 33.9 \pm 1.4 $ \\
Crowdflower progressive opinion & $ 70.6 $ & $ 71.6 \pm 0.5 $ & $ 75.9 \pm 1.3 $ & $ 71.7 \pm 1.1 $ \\
Crowdflower progressive stance & $ 54.3 $ & $ 63.8 \pm 1.0 $ & $ 61.5 \pm 1.3 $ & $ 60.6 \pm 1.4 $ \\
Crowdflower US economic performance & $ 75.6 $ & $ 75.3 \pm 0.1 $ & $ 76.5 \pm 0.4 $ & $ 77.3 \pm 0.1 $ \\
Customer complaint database & $ 54.5 $ & $ 55.9 \pm 0.1 $ & $ 56.4 \pm 0.1 $ & $55.4 \pm 0.1$ \\
News aggregator dataset & $ 95.2 $ & $ 96.3 \pm 0.0 $ & $ 96.5 \pm 0.0 $ & $ 96.2 \pm 0.0 $ \\
SMS spam collection & $ 98.5 $ & $ 99.3 \pm 0.2 $ & $ 99.3 \pm 0.2 $ & $ 95.1 \pm 2.2 $ \\
\midrule
Average & $72.7$ & $73.7$ & $74.0$ & $73.3$  \\
\midrule
Total number of params & \textemdash & $17\times$ & $9.9\times$ & $1.19\times$   \\
Trained params/task & \textemdash & $100\%$ & $52.9\%$ & $1.14\%$ \\
\bottomrule
\end{tabular}
\end{adjustbox}
\caption{
Test accuracy for additional classification tasks.
In these experiments we transfer from the BERT\textsubscript{BASE} model.
For each task and algorithm, the model with the best validation set accuracy is chosen.
We report the mean test accuracy and s.e.m. across runs with different random seeds.}
\label{tab:hub_results}
\vskip-3mm
\end{table*}
\renewcommand{\arraystretch}{1.0}

To further validate that adapters yields compact, performant, models, we test on additional, publicly available, text classification tasks.
This suite contains a diverse set of tasks:
The number of training examples ranges from $900$ to $330$k,
the number of classes ranges from $2$ to $157$,
and the average text length ranging from $57$ to $1.9$k characters.
Statistics and references for all of the datasets are in the appendix.

For these datasets, we use a batch size of $32$.
The datasets are diverse, so we sweep a wide range of learning rates:
$\{1 \cdot 10^{-5}, 3 \cdot 10^{-5}, 1 \cdot 10^{-4}, 3 \cdot 10^{-3}\}$.
Due to the large number of datasets, we select the number of training epochs from the set $\{20, 50, 100\}$ manually, from inspection of the validation set learning curves.
We select the optimal values for both fine-tuning and adapters; the exact values are in the appendix.

We test adapters sizes in $\{2, 4, 8, 16, 32, 64\}$.
Since some of the datasets are small, fine-tuning the entire network may be sub-optimal.
Therefore, we run an additional baseline: variable fine-tuning.
For this, we fine-tune only the top $n$ layers, and freeze the remainder.
We sweep $n\in\{1,2,3,5,7,9,11,12\}$.
In these experiments, we use the BERT\textsubscript{BASE} model with $12$ layers,
therefore, variable fine-tuning subsumes full fine-tuning when $n=12$.

\begin{figure*}[t]
\centering
\vskip-1mm
\begin{tabular}{cc}
GLUE (BERT\textsubscript{LARGE}) & Additional Tasks (BERT\textsubscript{BASE}) \\
\includegraphics[width=0.45\textwidth]{figures/glue_results_b.pdf}&
\includegraphics[width=0.45\textwidth]{figures/extra_tasks_plot.pdf}
\end{tabular}
\vskip-2mm
\caption{
Accuracy versus the number of trained parameters, aggregated across tasks.
We compare adapters of different sizes (orange) with fine-tuning the top $n$ layers, for varying $n$ (blue).
The lines and shaded areas indicate the $20$th, $50$th, and $80$th percentiles across tasks.
For each task and algorithm, the best model is selected for each point along the curve.
For GLUE, the validation set accuracy is reported.
For the additional tasks, we report the test-set accuracies.
To remove the intra-task variance in scores,
we normalize the scores for each model and task by subtracting the performance of full fine-tuning on the corresponding task.
}
\label{fig:tradeoff_alltasks}
\end{figure*}

\begin{figure*}[h!]
\centering
\vskip-1mm
\begin{tabular}{cc}
MNLI\textsubscript{m}(BERT\textsubscript{BASE}) & CoLA (BERT\textsubscript{BASE}) \\
\includegraphics[width=0.45\linewidth]{figures/param_efficiency_mnli.pdf}&
\includegraphics[width=0.45\linewidth]{figures/param_efficiency_cola.pdf}
\end{tabular}
\caption{
Validation set accuracy versus number of trained parameters for three methods:
(i) Adapter tuning with an adapter sizes $2^n$ for $n=0 \ldots 9$ (orange).
(ii) Fine-tuning the top $k$ layers for $k=1\ldots 12$ (blue).
(iii) Tuning the layer normalization parameters only (green).
Error bars indicate $\pm 1$ s.e.m. across three random seeds.}
\label{fig:tradeoff_glue}
\vskip-3mm
\end{figure*}

Unlike the GLUE tasks, there is no comprehensive set of state-of-the-art numbers for this suite of tasks.
Therefore, to confirm that our BERT-based models are competitive, we collect our own benchmark performances.
For this, we run a large-scale hyperparameter search over standard network topologies.
Specifically, we run the single-task Neural AutoML algorithm, similar to~\citet{zoph2017,wong2018transferautoml}.
This algorithm searches over a space of feedforward and convolutional networks,
stacked on pre-trained text embeddings modules publicly available via TensorFlow Hub\footnote{\url{https://www.tensorflow.org/hub}}.
The embeddings coming from the TensorFlow Hub modules may be frozen or fine-tuned.
The full search space is described in the appendix.
For each task, we run AutoML for one week on CPUs, using $30$ machines.
In this time the algorithm explores over $10$k models on average per task.
We select the best final model for each task according to validation set accuracy.

The results for the AutoML benchmark (``no BERT baseline''), fine-tuning, variable fine-tuning, and adapter-tuning are reported in Table~\ref{tab:hub_results}.
The AutoML baseline demonstrates that the BERT models are competitive.
This baseline explores thousands of models, yet the BERT models perform better on average.
We see similar pattern of results to GLUE.
The performance of adapter-tuning is close to full fine-tuning ($0.4\%$ behind).
Fine-tuning requires $17\times$ the number of parameters to BERT\textsubscript{BASE} to solve all tasks.
Variable fine-tuning performs slightly better than fine-tuning, whilst training fewer layers.
The optimal setting of variable fine-tuning results in training $52\%$ of the network on average per task, reducing the total to $9.9\times$ parameters.
Adapters, however, offer a much more compact model.
They introduce $1.14\%$ new parameters per task, resulting in $1.19\times$ parameters for all $17$ tasks.

\subsection{Parameter/Performance trade-off\label{sec:param_efficiency}}

The adapter size controls the parameter efficiency, smaller adapters introduce fewer parameters, at a possible cost to performance.
To explore this trade-off, we consider different adapter sizes, and compare to two baselines:
(i) Fine-tuning of only the top $k$ layers of BERT\textsubscript{BASE}.
(ii) Tuning only the layer normalization parameters.
The learning rate is tuned using the range presented in Section~\ref{sec:glue}.

Figure~\ref{fig:tradeoff_alltasks} shows the parameter/performance trade-off aggregated over all classification tasks in each suite (GLUE and ``additional'').
On GLUE, performance decreases dramatically when fewer layers are fine-tuned.
Some of the additional tasks benefit from training fewer layers, so performance of fine-tuning decays much less.
In both cases, adapters yield good performance across a range of sizes two orders of magnitude fewer than fine-tuning.

Figure~\ref{fig:tradeoff_glue} shows more details for two GLUE tasks: MNLI\textsubscript{m} and CoLA.
Tuning the top layers trains more task-specific parameters for all $k>2$.
When fine-tuning using a comparable number of task-specific parameters, the performance decreases substantially compared to adapters.
For instance, fine-tuning just the top layer yields approximately $9$M trainable parameters and $77.8 \% \pm 0.1 \%$ validation accuracy on MNLI\textsubscript{m}.
In contrast, adapter tuning with size $64$ yields approximately $2$M trainable parameters and $83.7\% \pm 0.1 \%$  validation accuracy.
For comparison, full fine-tuning attains $84.4 \% \pm 0.02 \%$ on MNLI\textsubscript{m}.
We observe a similar trend on CoLA.

As a further comparison, we tune the parameters of layer normalization alone.
These layers only contain point-wise additions and multiplications, so introduce very few trainable parameters: $40$k for BERT\textsubscript{BASE}.
However this strategy performs poorly: performance decreases by approximately $3.5\%$ on CoLA and $4\%$ on MNLI.

To summarize, adapter tuning is highly parameter-efficient, and produces a compact model with a strong performance, comparable to full fine-tuning.
Training adapters with sizes $0.5-5\%$ of the original model,
performance is within $1\%$ of the competitive published results on BERT\textsubscript{LARGE}.

\subsection{SQuAD Extractive Question Answering}

\begin{figure}[t]
\centering
\vskip-2mm
\includegraphics[width=0.9\linewidth]{figures/squad_adapters_baseline.pdf}
\caption{
Validation accuracy versus the number of trained parameters for SQuAD v1.1.
Error bars indicate the s.e.m. across three seeds, using the best hyperparameters.
}
\label{fig:squad}
\vskip-5mm
\end{figure}

Finally, we confirm that adapters work on tasks other than classification by running on SQuAD v1.1~\citep{rajpurkar2018}.
Given a question and Wikipedia paragraph, this task requires selecting the answer span to the question from the paragraph.
Figure~\ref{fig:squad} displays the parameter/performance trade-off of fine-tuning and adapters on the SQuAD validation set.
For fine-tuning, we sweep the number of trained layers, learning rate in $\{3\cdot 10^{-5}, 5\cdot 10^{-5}, 1\cdot 10^{-4}\}$, and number of epochs in $\{2,3,5\}$.
For adapters, we sweep the adapter size, learning rate in $\{3\cdot 10^{-5}, 1\cdot 10^{-4}, 3\cdot 10^{-4}, 1\cdot 10^{-3}\}$, and number of epochs in $\{3,10,20\}$.
As for classification, adapters attain performance comparable to full fine-tuning, while training many fewer parameters.
Adapters of size $64$ ($2\%$ parameters) attain a best F1 of $90.4$\%, while fine-tuning attains $90.7$.
SQuAD performs well even with very small adapters, those of size $2$ ($0.1\%$ parameters) attain an F1 of $89.9$.

\subsection{Analysis and Discussion\label{sec:discussion}}

We perform an ablation to determine which adapters are influential.
For this, we remove some trained adapters and re-evaluate the model (without re-training) on the validation set.
Figure~\ref{fig:ablation_and_init} shows the change in performance when removing adapters from all continuous layer spans.
The experiment is performed on BERT\textsubscript{BASE} with adapter size $64$ on MNLI and CoLA.

First, we observe that removing any single layer's adapters has only a small impact on performance.
The elements on the heatmaps' diagonals show the performances of removing adapters from single layers, where largest performance drop is $2\%$.
In contrast, when all of the adapters are removed from the network,
the performance drops substantially:
to $37\%$ on MNLI and $69\%$ on CoLA -- scores attained by predicting the majority class.
This indicates that although each adapter has a small influence on the overall network, the overall effect is large.

Second, Figure~\ref{fig:ablation_and_init} suggests that adapters on the lower layers have a smaller impact than the higher-layers.
Removing the adapters from the layers $0-4$ on MNLI barely affects performance.
This indicates that adapters perform well because they automatically prioritize higher layers.
Indeed, focusing on the upper layers is a popular strategy in fine-tuning~\citep{howard2018universal}.
One intuition is that the lower layers extract lower-level features that are shared among tasks, while the
higher layers build features that are unique to different tasks.
This relates to our observation that for some tasks, fine-tuning only the top layers outperforms full fine-tuning, see Table~\ref{tab:hub_results}.

Next, we investigate the robustness of the adapter modules to the number of neurons and initialization scale.
In our main experiments the weights in the adapter module were drawn from
a zero-mean Gaussian with standard deviation $10^{-2}$, truncated to two standard deviations.
To analyze the impact of the initialization scale on the performance, we test standard deviations in the interval $[10^{-7},1]$.
Figure~\ref{fig:ablation_and_init} summarizes the results.
We observe that on both datasets, the performance of adapters is robust for standard deviations below $10^{-2}$.
However, when the initialization is too large, performance degrades, more substantially on CoLA.

To investigate robustness of adapters to the number of neurons, we re-examine the experimental data from Section~\ref{sec:glue}.
We find that the quality of the model across adapter sizes is stable,
and a fixed adapter size across all the tasks could be used with small detriment to performance.
For each adapter size we calculate the mean validation accuracy across the eight
classification tasks by selecting the optimal learning rate and number of epochs\footnote{
We treat here MNLI\textsubscript{m} and MNLI\textsubscript{mm} as separate tasks.
For consistency, for all datasets we use accuracy metric and exclude the regression STS-B task.}.
For adapter sizes $8$, $64$, and $256$, the mean validation accuracies are $86.2\%$, $85.8\%$ and $85.7\%$, respectively.
This message is further corroborated by Figures~\ref{fig:tradeoff_glue} and~\ref{fig:squad},
which show a stable performance across a few orders of magnitude.

Finally, we tried a number of extensions to the adapter's architecture
that did not yield a significant boost in performance.
We document them here for completeness.
We experimented with
(i) adding a batch/layer normalization to the adapter,
(ii) increasing the number of layers per adapter,
(iii) different activation functions, such as tanh,
(iv) inserting adapters only inside the attention layer,
(v) adding adapters in parallel to the main layers, and possibly with a multiplicative interaction.
In all cases we observed the resulting performance to be similar to the bottleneck proposed in Section~\ref{sec:bottleneckadapter}.
Therefore, due to its simplicity and strong performance, we recommend the original adapter architecture.

\begin{figure*}
\centering
\begin{tabular}[t]{ccc}
MNLI\textsubscript{m} & CoLA & \\
\includegraphics[height=3.7cm]{figures/mnli_ablation_heatmap.pdf}&
\includegraphics[height=3.7cm]{figures/cola_ablation_heatmap.pdf}&
\raisebox{-0.2cm}{\includegraphics[height=3.9cm]{figures/init_study.pdf}}
\end{tabular}
\caption{
\textbf{Left, Center:}
Ablation of trained adapters from continuous layer spans.
The heatmap shows the relative decrease in validation accuracy to the fully trained adapted model.
The \emph{y} and \emph{x} axes indicate the first and last layers ablated (inclusive), respectively.
The diagonal cells, highlighted in green, indicate ablation of a single layer's adapters.
The cell in the top-right indicates ablation of all adapters.
Cells in the lower triangle are meaningless, and are set to $0\%$, the best possible relative performance.
\textbf{Right:}
Performance of BERT\textsubscript{BASE} using adapters with different initial weight magnitudes.
The \emph{x}-axis is the standard deviation of the initialization distribution.
}
\label{fig:ablation_and_init}
\vskip-4mm
\end{figure*}
